\section{Conclusion and future work}
\label{sec:conclusion}

\slowlab{} turns experimental design into a sequential decision problem with physical
capacity, delayed terminal outcomes, correlated measurements, economic costs and
irreversible commitments. Version~2 adds a distinction that the original environment
missed: a slow experiment can still be inspected while it runs. An agent can use
measurements available at time $t$, while future yield, future profit and latent simulator
state remain hidden and committed treatments remain fixed.

The Phase~2 audit changes the evidential status of the original study. The finite-atom
estimator did not support the claimed design--decision decomposition, component fields
created an information bypass, and normalised density was used as a physical plant density.
Because the resulting observations affected later choices, the old closed-loop trajectories
cannot be repaired by rescoring. Claims about saturation, independent design and reader
failures, tool effects, transfer gains and first-cycle deterioration are therefore
withdrawn. The archived runs remain useful for interface and executed-design diagnostics.

The replacement estimand conditions every comparison on the same visible history. A
continuous crop--economics posterior passes the one-factor Sanity calibration and supports
held-out estimates of Bayes risk and recommendation excess risk, with an explicit lower-bound
flag when the best action is searched over a finite set. Its present calibration does not
justify precise efficiency claims on the higher-dimensional tasks.

\paragraph{Future work.} The next experiments should first regenerate closed-loop data under
the corrected density and observation model. They should then cross designs and readers,
separate design and inference tools, and compare terminal-only feedback with scheduled and
agent-selected interim inspection at matched resource and measurement budgets. Only effects
that replicate on held-out sites should return to the abstract and conclusion.
